\section{The Integer Ladder}

The base of the theory has no real numbers and no decimal numbers anywhere. This is not a matter of taste. It is the first axiom, axiom A1, and it is the deepest constraint the whole framework obeys. Everything that lives at the base is either an integer or one of three tone values. The reals are admitted later, but only as things we measure after the discrete machine has run, never as anything the machine stores or steps on.

This section makes the commitment precise, shows why it is not arbitrary, and traces the single thread that runs from the bare integers all the way up to the substrate the rest of the paper builds on. The thread is short and every rung rests on the one beneath it. Nowhere along it is there a real-valued dial to turn.

\begin{principle}[The base must be exactly stored and exactly stepped]
The base must be a thing that could be written down without rounding and advanced without rounding. It is a finite, exact object at every beat, not a continuum.
\end{principle}

\subsection{The discreteness axiom}

We name the constraint and then unpack it.

\begin{axiom}[A1, Discrete]
There are no real numbers and no decimal numbers anywhere in the base. The mesh holds a finite count of docks at any beat. Each dock has a finite count of sites. Each site carries a tone drawn from the three-valued alphabet $\{-1, 0, +1\}$. Time is an integer count of beats. The symmetry group of the substrate is finite, the group $F_4$ of order $1152$. No quantity in the base is a real number, and no step needs one.
\end{axiom}

Every word of A1 carries weight. The mesh is the whole $\{3,4,3,4\}$ honeycomb, a weave of docks, finite at any beat and ever growing along its wake. A dock is one cell, a here and now, and it has exactly $24$ sites, one for each direction out of the cell. There is no twenty-fifth slot, no rest site, no still home. A site holds one tone, and the tone is one of three values, read as pain, peace, and pleasure. The lean is the value order $-1 < 0 < +1$ that turns those tones into signed charge. The knit is the law that changes the state each beat, collide then stream. The wake is the law that grows the space each beat. The beat is the discrete tick of time. None of these is a real number, and none of them needs one.

Breaking A1 is a cheat, and the theory treats it as one. A stored real coordinate is forbidden. A real-valued coupling constant is forbidden. A continuous time step is forbidden. The base is built from the integers $\mathbb{Z}$ and from the three-valued alphabet $\{-1, 0, +1\}$, and from nothing else.

The deepest consequence of A1, taken together with the emergence gap, is the two-layer structure of the whole theory. The base is discrete, call it Layer 0. The self and every recognizable piece of physics are emergent, Layer 1 and above, the coarse-grained average of enormously many discrete grains. The emergent layer may look perfectly continuous even though every grain underneath it is a single tone in $\{-1, 0, +1\}$. Discrete below, smooth above. The two never contradict, because they live at different scales.

\subsection{The ladder, each rung resting on the one below}

A fair worry about any cellular model is that the choice of cell looks arbitrary. Why a honeycomb, why this honeycomb, why this knit. The answer is that the whole substrate rests on the least arbitrary thing there is, the integers. The reflection groups, the crystals, and the addressing are not separate inventions. They are different views of one object built from $\mathbb{Z}$. There are no real numbers and no tunable dials anywhere in the foundation. The tower runs as follows.

\begin{description}
\item[Rung 0, the integers.] Start with $\mathbb{Z}$. Nothing is more canonical. No parameters, no reals, no randomness.
\item[Rung 1, integer symmetries.] Ask for the symmetries generated by reflections, where the only data is a set of integers saying how the mirrors meet. Mirror $i$ and mirror $j$ meet at an angle $\pi / m_{ij}$ for an integer $m_{ij}$.
\item[Rung 2, Coxeter groups.] A group built from reflections whose meeting angles are fixed by integers in exactly this way is a Coxeter group. Its entire data is a small integer matrix, the Coxeter diagram. Still nothing but integers.
\item[Rung 3, tessellations.] Let a Coxeter group act on a space by its reflections. Reflect one fundamental chamber over and over, and the orbit fills the space with copies. That filling is a tessellation. The algebraic group and the geometric tiling are the same thing seen two ways.
\item[Rung 4, the specific crystals.] Choose the integers in the diagram and you get one specific tiling, named by its Schl\"afli symbol, such as $\{7,3\}$, $\{5,3,4\}$, or the substrate of this paper, $\{3,4,3,4\}$. These are not separate constructions. They are Rung 3 with particular integer settings. One machine on different settings.
\item[Rung 5, the substrate.] On the chosen tessellation the vibes live. Tones sit on the sites of the docks, and the knit runs along the directions. The substrate inherits, for free, everything the lower rungs supply.
\item[Rung 6, the generating automaton.] Every hyperbolic Coxeter group is an automatic group. A finite integer state machine generates the whole infinite tessellation deterministically, with no randomness. For $\{3,4,3,4\}$ this is the solved neighbour automaton together with the $O(\log n)$ addressing. The substrate generates and navigates itself, starting from $\mathbb{Z}$, by a finite rule.
\end{description}

\begin{result}[The substrate is one integer setting of a canonical machine]
The whole substrate, from its symmetry group to its tiling to its self-generating automaton, is obtained by choosing integers in a Coxeter diagram. There is no free real-number dial anywhere in the foundation. The choice of $\{3,4,3,4\}$ is the choice of which small integers go in the diagram, nothing more.
\end{result}

\subsection{One object, many views}

Because the rungs are different views of one thing, the substrate wears many faces at once. For $\{3,4,3,4\}$ the Coxeter system $[3,4,3,4]$ is, simultaneously, all of the following.

\begin{table}[ht]
\centering
\begin{tabular}{@{}ll@{}}
\toprule
view & what it is \\
\midrule
algebraic group & five reflections \\
geometric honeycomb & $24$-cells filling $\mathbb{H}^4$ \\
combinatorial addressing tree & the splitting matrix, Zeckendorf-style digits \\
dynamical generator & the neighbour automaton, beat by beat \\
growth law & the shell recurrence, with its algebraic Perron root as the central constant \\
\bottomrule
\end{tabular}
\caption{One object, several faces, for the substrate $\{3,4,3,4\}$.}
\label{tab:many-views}
\end{table}

These are not analogies and not loose correspondences. They are literally the same object looked at in several ways. The algebra and the geometry and the dynamics all read off the same integer data.

\subsection{Exact integer and ternary arithmetic}

Every operation in the base is exact. Nothing rounds and nothing drifts.

\begin{table}[ht]
\centering
\begin{tabular}{@{}lll@{}}
\toprule
quantity & type & exact \\
\midrule
a tone, the value in one site & ternary $\{-1, 0, +1\}$ & yes \\
a dock's net charge & bounded integer, the lean-sum of $24$ trits & yes \\
total charge over the mesh & integer, conserved & yes \\
the dock count & integer, finite per beat & yes \\
the site count, the degree & integer, $24$ & yes \\
time & integer count of beats & yes \\
a neighbour index & integer & yes \\
the symmetry group order & integer, $1152$ & yes \\
\bottomrule
\end{tabular}
\caption{Every base quantity is an integer or a tone, and every operation on it is exact.}
\label{tab:exact-arithmetic}
\end{table}

The collide phase of the knit is a permutation of a finite set of dock states. A permutation is exact, with no rounding anywhere in it. The stream phase of the knit is a relabeling of integer site indices, moving each tone one dock along its direction. A relabeling is exact too. Charge conservation is integer addition that comes out identical every beat, by construction of the collide table.

\begin{proposition}[No float in a beat]
A single beat is a permutation of a finite state set followed by a relabeling of integer indices. Both are exact operations on integers and tones. There is no real-valued update at any point in a beat, so nothing accumulates error. The same start state run for a billion beats is bit-for-bit reproducible.
\end{proposition}

\subsection{Where real numbers are allowed}

Real numbers are not banished from the theory. They are banished from the base. They are welcome in the measurement, the layer where we coarse-grain the discrete machine and read quantities off it. After the integer machine has run, the things we fit and observe may be real.

\begin{table}[ht]
\centering
\begin{tabular}{@{}lll@{}}
\toprule
real-valued thing & where it lives & base or output \\
\midrule
a tone, a charge, a site value & the base buffer & base, always integer or ternary \\
an emergent coordinate & a measured position on the horosphere & output \\
an eigenvalue of a measured operator & a fitted dispersion or spectrum & output \\
an emergent dimension & a scaling exponent we fit & output \\
a continuous symmetry, Lorentz, $1/r$ attraction & the coarse-grained infrared & output, emergent \\
\bottomrule
\end{tabular}
\caption{Reals appear only as measured outputs, never in the base.}
\label{tab:reals-as-output}
\end{table}

So a dispersion of the form $\cos E = \cos m \cos k$, the Perron root near $18.3$ that governs growth, the emergent spatial dimension, all of these are quantities we read off the discrete machine. They are never stored inside it and never used to step it. The arrow of time and the attraction are emergent in the same sense. The lean orders the tones, the wake feeds in low entropy at the frontier, and the direction of time appears as a coarse-grained output of that asymmetry, not as a base ingredient.

\subsection{Why this matters}

The integer commitment is not bookkeeping. It buys four properties that the theory leans on heavily.

\begin{description}
\item[Exactly storable.] The whole universe at any beat is one flat array of small integers, every entry a tone in $\{-1, 0, +1\}$. A finite, exact object that could in principle be written down in full.
\item[Exactly steppable.] The beat is a permutation followed by a relabeling. Both are exact. There is no real-valued update, no rounding, and no drift.
\item[Reproducible.] Determinism, the axiom A2, together with integer arithmetic, means the same start gives the same run, bit-for-bit, forever. To measure at a larger scale we grow the size of the mesh, never a seed and never a shrunk real parameter.
\item[No free parameter.] To pick the base is only to pick which integers go in the Coxeter diagram. The object beneath all the views is one integer machine. There is no free real-number dial anywhere in the foundation.
\end{description}

So $\{3,4,3,4\}$ is not an arbitrary choice. It is one integer setting of the canonical Coxeter machine, and which setting is forced is settled by the selection criteria taken up in the tessellation survey and the dimension ladder. The choice of substrate reduces, in the end, to a choice among small integers, and the selection argument is what fixes them.

\begin{principle}[The integers carry the laws]
The integers are the bones of the theory. The small ones carry the laws. The base stores only integers and tones, steps only by exact permutation and relabeling, and admits real numbers solely as the coarse-grained things we measure once the discrete machine has run.
\end{principle}
